
\documentclass[12pt]{article}

\usepackage[margin=1in]{geometry}
\usepackage{amsmath, amssymb}
\usepackage{siunitx}
\usepackage{booktabs}
\usepackage{longtable}
\usepackage{tabularx}
\usepackage{array}
\usepackage{enumitem}
\usepackage{hyperref}
\usepackage{listings}
\usepackage{xcolor}

\hypersetup{
    colorlinks=true,
    linkcolor=blue,
    urlcolor=blue,
    citecolor=blue
}

\lstset{
    basicstyle=\ttfamily\small,
    breaklines=true,
    columns=fullflexible,
    frame=single
}

\title{Ionizing and Non-Ionizing Radiation with Emphasis on X-Rays}
\author{Zachary Tullier}
\date{6/29/26}

\begin{document}

\maketitle
\section{Radiation Dose (Dosage), Dosimetry, and Dosimetry Methods}

\subsection{Radiation Dose}
Radiation dose is a quantitative measure of the energy deposited by ionizing radiation in matter, particularly biological tissues.

\subsubsection{Exposure}
\begin{equation}
X=\frac{dQ}{dm}
\end{equation}
where $Q$ is electric charge and $m$ is the mass of air. The SI unit is C/kg, with
\[
1~\mathrm{R}=2.58\times10^{-4}\,\mathrm{C/kg}.
\]

\subsubsection{Absorbed Dose}
\begin{equation}
D=\frac{dE}{dm}
\end{equation}
The SI unit is the Gray:
\[
1~\mathrm{Gy}=1~\mathrm{J/kg}=100~\mathrm{rad}.
\]

\subsubsection{Equivalent Dose}
\begin{equation}
H_T=Dw_R
\end{equation}

\begin{table}[h]
\centering
\begin{tabular}{lc}
\hline
Radiation & $w_R$\\ \hline
X-rays & 1\\
Gamma rays & 1\\
Beta particles & 1\\
Protons & 2\\
Alpha particles & 20\\
Neutrons & 5--20\\
\hline
\end{tabular}
\caption{Radiation weighting factors.}
\end{table}

\subsubsection{Effective Dose}
\begin{equation}
E=\sum_T w_T H_T
\end{equation}

\section{Radiation Dosimetry}
Radiation dosimetry is the science of measuring absorbed radiation dose.

Applications include radiation therapy, diagnostic radiology, nuclear medicine, industrial radiography, environmental monitoring, occupational protection, and space radiation monitoring.

\section{Types of Dosimeters}

\subsection{Film Badge}
Film badges darken in proportion to radiation exposure due to silver bromide ionization.

\begin{verbatim}
Radiation
   |
Film
   |
Latent Image
   |
Development
   |
Dose
\end{verbatim}

Advantages: inexpensive, permanent record.
Disadvantages: single-use, affected by heat and humidity.

\subsection{Thermoluminescent Dosimeter (TLD)}
Heating releases trapped electrons from LiF, CaSO$_4$, or CaF$_2$ crystals.

\[
L\propto D
\]

\subsection{Optically Stimulated Luminescence (OSL)}
Laser light releases trapped electrons from Al$_2$O$_3$:C crystals. The emitted light is proportional to dose.

\subsection{Ionization Chamber}
\begin{equation}
I=\frac{Q}{t}
\end{equation}

Gas ionization current is measured with an electrometer.

\subsection{Geiger--M\"uller Counter}
Radiation initiates a Townsend avalanche producing detectable electrical pulses.

\subsection{Semiconductor Dosimeters}
Radiation creates electron-hole pairs in silicon or diamond.

\[
Q\propto D
\]

\subsection{Chemical Dosimeters}
Fricke dosimeters measure:
\[
\mathrm{Fe^{2+}\rightarrow Fe^{3+}}
\]

\subsection{Alanine Dosimeters}
Alanine crystals form stable free radicals measured by ESR spectroscopy.

\section{Comparison of Dosimetry Methods}

\begin{table}[h]
\centering
\begin{tabular}{llll}
\hline
Dosimeter & Principle & Readout & Use\\ \hline
Film Badge & Film darkening & Delayed & Personnel\\
TLD & Thermoluminescence & Delayed & Medical\\
OSL & Optical luminescence & Delayed & Personnel\\
Ionization Chamber & Gas ionization & Real time & Calibration\\
GM Counter & Gas avalanche & Real time & Surveys\\
Semiconductor & Electron-hole pairs & Real time & Radiotherapy\\
Fricke & Chemical oxidation & Delayed & High dose\\
Alanine & ESR & Delayed & Industrial\\
\hline
\end{tabular}
\caption{Comparison of dosimetry methods.}
\end{table}

\section{Applications of Dosimetry}
\begin{itemize}
\item Diagnostic radiology
\item Radiation therapy
\item Nuclear medicine
\item Occupational monitoring
\item Environmental monitoring
\item Industrial irradiation
\item Space radiation monitoring
\end{itemize}  
\end{document}